\documentclass[12pt, a4paper]{report}
\edef\restoreparindent{\parindent=\the\parindent\relax}
\usepackage[UKenglish]{babel}
\usepackage[bibstyle=ieee, dashed=false, sorting=nty]{biblatex}
\usepackage[labelfont=bf]{caption}
\usepackage{colortbl}
\usepackage{csquotes}
\usepackage{fancyhdr}
\usepackage{float}
\usepackage[bottom]{footmisc}
\usepackage{graphicx}
\usepackage[hidelinks]{hyperref}
\usepackage{parskip}
\usepackage{pgfgantt}

\linespread{1.2}
\restoreparindent

\pagestyle{fancy}
\fancyhf{}
\fancyhead[C]{\leftmark}
\fancyfoot[C]{\thepage}

\addbibresource{references.bib}

\begin{document}
\begin{titlepage}
	\centering
	\includegraphics[width=10cm]{tuos_logo}\par\vspace{1cm}
	\vspace{1cm}

	{\huge\bfseries Finding Security Issues in (Open Source) Software Repositories\par}
	\vspace{1cm}

	{\Large Zer Jun Eng\par}
	\vspace{1cm}

	supervised by\par Dr.~Achim \textsc{Brucker}
	\vfill

	{This report is submitted in partial fulfilment of the requirement for the degree of MEng Software
		Enginnering by Zer Jun Eng}
	\vfill

	{\large COM3610}
	\vfill

	{\large \today}
\end{titlepage}

\pagenumbering{roman}

\chapter*{Declaration}
All sentences or passages quoted in this report from other people's work have been specifically
acknowledged by clear cross-referencing to author, work and page(s). Any illustrations that are not
the work of the author of this report have been used with the explicit permission of the originator
and are specifically acknowledged. I understand that failure to do this amounts to plagiarism and
will be considered grounds for failure in this project and the degree examination as a whole.
\vspace{2cm}

\noindent \begin{tabular}{llp{4.5cm}}
	Name & : & Zer Jun Eng \\ \cline{3-3}
	\\ [-0.5em]
	Date & : & \today      \\ \cline{3-3}
\end{tabular}

\newpage

% \chapter*{Abstract}

\chapter*{Acknowledgements}
I would firstly like to thank my parents for their unconditional love and the full financial supoprt
throughout my university life. It would not be possible for me to finish this project and my course
without them.

I would also like to thank my supervisor, Dr. Achim Brucker, who are continuously providing
constructive advice for my project. I am honoured to work with you, and I look forward to working
with you in the future.

\newpage

\tableofcontents

% \listoffigures

\listoftables

\newpage

\pagenumbering{arabic}

\chapter{Introduction}
\section{Background}
Free/Libre and Open Source Software (\textbf{FLOSS}) is a type of software which its license allows
the users to inspect, use, modify and redistribute the software's source code \cite{crowston_2012}.
Since the introduction of version control system, many repositories hosting site such as
SourceForge\cite{sourceforge}, Google Code \cite{google_code}, and GitHub \cite{github} have been
launched. As a result, the participation of global communities into FLOSS projects has started to
grow and different contributions were made to improve the softwares quality, which included fixing
the software vulnerabilities \cite{dabbish_2012}.

Building a secure software is expensive, difficult, and time-consuming. It is necessary to know when
and how a security vulnerability is fixed throughout the software lifespan. One of the factors is
that many software components are commonly reused in both open source and closed source softwares
\cite{khan_2001}, and a lot of information is required to decide whether or not to perform the
security update provided by the vendor of the components. Therefore, having a list of changelogs or
informative Git commit messages that record the fixed security vulnerabilities is helpful. However,
some developers believe that public disclosure of security vulnerabilities patch is dangerous, thus
vulnerability fixing commits are not commonly identified in some open source software repositories
to prevent malicious exploits \cite{arora_2005}. To address these issues, a repository mining tool
that investigate commit messages and identify vulnerable software components can be developed to
reduce the time and cost required to mitigate the vulnerabilities. Moreover, the mining tool should
be applicable to all types of software projects that are using Git as their version control system.

\section{Objectives} \label{sec:objectives}
\begin{itemize}
	\item Identify the security patterns of the most popular security issues in OWASP Top Ten Project.
	The patterns should be expressed using regular expressions.
	\item Develop a repository mining tool to search through the commit history of a repository and
	find a list of commit messages that match the patterns. The list should be produced in a suitable
	file format such as JSON, XML, or CSV.
	\item Extend the mining tool which checks the code difference in the commits found to obtain the
	actual commits fixing the security vulnerabilities. This extension should separate from the mining
	process to make the mining results easier to verify and debug.
\end{itemize}

\section{Challenges} \label{sec:challenges}
\begin{itemize}
	\item \textbf{Data}: There are a large number of open source repositories available on GitHub.
	However, it is challenging to find a set of sample repositories that can produce accurate and
	consistent results.
	\item \textbf{Misclassification}: Commit messages for a same vulnerability patch are not always
	the same, thus misclassification of commit messages is inevitable. Using regular expressions to
	match the patterns in the mining process does not guarantee correctness of the result.
	\item \textbf{Evaluation}: After mining a list of commits that contain the identified patterns in
	its message, the evaluation process might not correctly locate the lines of code that addressed
	the security vulnerability. It might be required to perform manual evaluation to correctly
	identify some of the results.
	\item \textbf{Time}: Large repository such as Linux which has more than 780,000 commits in total
	\cite{linux_repo} could be extremely time-consuming for the repository mining tool to complete the
	search and evaluation process.
\end{itemize}

\section{Report Structure}
\textbf{Chapter 2} reviews a range of academic articles, theories and previous studies that
is related to this project, as well as investigating the techniques and tools to be used.

\noindent\textbf{Chapter 3} is a list of detailed requirements and a thorough analysis for design,
implementation and testing stage. Some core decisions are reviewed in the analysis part to ensure
the feasibility of the project.

\noindent\textbf{Chapter 4} is a comparison between different design concepts, where the advantages
and disadvantages of difference approach are stated. The chosen design is justified with suitable
diagrams provided including wireframes and UML.

\noindent\textbf{Chapter 5} describes the implementation process by highlighting novel aspects to
the algorithms used. Testing are performed by following a suitable model to evaluate the
implementation.

\noindent\textbf{Chapter 6} presents all the results along with critical discussions about the main
findings,	and outlines the possible improvements that could be made in the future work.

\noindent\textbf{Chapter 7} summarises the main points of previous chapters and emphasise the
results found.

\section{Relationship to Degree Programme}
This project focuses on the reserach of real-world software security problems, and offers a valuable
insights into the computer security. By studying the patterns of security vulnerabilities patch,
practical knowledge for building and ensuring a secure system could be gained. Moreover, the
difficulty of improving software security could be experienced during the evaluation process in this
project. This relates to the Software Engineering degree as it requires a good understanding in
version control system and it aims to improve softwares quality by reducing the time and effort
needed to find security vulnerabilities in the source code.

\chapter{Literature Review}
This chapter will start with the background contents of the project, and then focus on discussing
the security aspect of open source softwares. Lastly, previous and existing relevant work are
reviewed and a critical analysis is provided for the comparison of these resources and this project.

\section{Open Source Security}
The security of open source softwares mostly rely on the collaboration of the community with the
developers. It is deduced that the power of open data and crowdsourcing will make open source
security more reliable \cite{hoepman_2007, witten_2001}, and provides more flexibility and freedom
over the security option to their users \cite{payne_2002}. However, when it comes to publishing the
vulnerability information, it is suggested that the list of unconfirmed vulnerabilities should not
be published publicly to protect the users from potential harms \cite{schryen_2011}.

Arora, Nandkumar and Telang \cite{arora_2006} have shown that vulnerabilities that are either secret
or published but not patched attract fewer attacks than patched vulnerabilities. This implies that
developers might secretly include a patch into some of the commits, and it is not explicitly
recorded in the commit messages. It is a rational approach for not disclosing the work attempted to
fix a vulnerability, but this could make it harder to trace back the solution if a similar
vulnerability is discovered in the future. Therefore, it would be very useful for the developers if
the mining tool developed in this project could detect the silent patches.

\section{Common Vulnerabilities and Exposures}
The Common Vulnerabilities and Exposures (\textbf{CVE}) is a project launched by Mitre Corporation
and sponsored by the National Cyber Security Division of the United States Department of Homeland
Security \cite{cve}. The CVE system provides the computer security community with a complete list of
publicly known security vulnerabilities, and each vulnerability is identified by a unique CVE ID
number. It is now the standardised solution and industry-recognised standard for identifying
vulnerabilities and exposures. However, developers and vendors are not required to publish security
vulnerabilities of their projects in CVE format, and they are allowed to use their own naming scheme
for the vulnerabilities, even though the same vulnerability has been recorded in the CVE list.

\section{Common Weakness Enumeration}
The Common Weakness Enumeration (\textbf{CWE}) is another project of Mitre Corporation \cite{cwe}
that organises the software weaknesses into a list of different categories, known as the CWE list.
Software weaknesses are defined as the errors that can lead to software vulnerabilities, which
includes buffer overflows, authentication errors, code injection, etc. \cite{cwe_faq}. The CWE is
now a formal standard for representing software weaknesses. Each entry in the CWE are identified by
a unique ID number, and contained detailed information about the specific weakness.

\section{Security Issues in Open Source Softwares}
The Open Web Application Security Project (\textbf{OWASP}) is a worldwide non-profit organization
committed to improve and raise the awareness of software security \cite{owasp_home}. The project
members of OWASP have worked together to produce a list of the most critical web application
security risks based on the community feedback and comprehensive data contributed by different
organizations. The list consists of ten categories of security attacks which are considered to be
the most dangerous and popular in the recent years. One of the vulnerabilities that is closely
related to this project is \textit{Using Components with Known Vulnerabilities}, which will be
extensively discussed. In addition, \textit{Injection} and \textit{Broken Authentication} will be
discussed as a reference for this project because they are listed as the top two vulnerabilities in
OWASP Top Ten 2017 \cite{owasp_top10}.

\subsection{Using Components with Known Vulnerabilities} \label{subsec:components}
Components such as plugins, libraries, and modules can often be found in different parts of a
software. Third-party components are increasingly being integrated into softwares to reduce the
amount of time and effort required for development \cite{balzarotti_2006}, but they also increase
the risk of vulnerabilities being introduced into the softwares. These components are mostly
maintened by different developers or organisations hence it is unable to guarantee that all the
components used will be provided with the latest security patch.

Cadariu et al. \cite{cadariu_2015} used OWASP Dependency Check tool \cite{owasp_dependency} to find
all known vulnerabilities in proprietary softwares written in Java. However, their results contained
a considerable amount of both false positives and false negatives, which directly affected the
precision of their technique. Therefore, it is also expected that the final results produced in this
project might encounter the same problem.

\subsection{Injection} \label{subsec:injection}
An injection attack is the exploitation of a software vulnerability where the attacker injects
malicious code into the software applications and perform harmful executions. The most common types
of injection attacks are Structured Query Language (\textbf{SQL}) injection, cross-site scripting,
and shell injection \cite{pietraszek_2006}. Among these attacks, SQL injection attacks required the
most effort to fix \cite{othmane_2015}. This is because a successful SQL injection attack could
delete an entire database that could cause a massive data loss in extreme situation. SQL injection
attacks are very widespread and easy to discover especially when attackers have access to the source
code, or they could use a code scanner tool such as sqlmap \cite{sqlmap} to find all possible ways
of the injection attacks.

Vast majority of web applications make use of JavaScript to add automation and interactivity to
their web pages to enhances the user experience. However, the growing usage of JavaScript has also
increased the frequency and popularity of cross-site scripting (\textbf{XSS}) attacks. There three
types of XSS attacks are: \textit{reflected XSS}, \textit{stored XSS}, and \textit{DOM XSS}
\cite{kieyzun_2009}. While the attack approach of different XSS methods may vary, their end goal is
to steal or access the sensitive information of the victims \cite{vogt_2007}. XSS vulnerabilities
could be easily detected and exploited at the same time, for example, the open source tool XSStrike
\cite{xsstrike} can scan a web page for possible {XSS vulnerabilities. A promising approach to
tackle this problem would be adding an extra security mechanism for handling the user inputs on the
web pages.

\subsection{Broken Authentication} \label{subsec:authentication}
Broken authentication happens when the attackers are allowed to perform malicious actions such as
brute force dictionary attack on the authentication system. It might also be one of the outcome of a
successful injection attack. According to Huluka and Popov \cite{huluka_2012}, this vulnerability
are very prevalent and has various causes, in which the lack of attention to security details is the
most critial because developers often overlook certain scenarios which are likely to be exploited by
attackers.

\section{Mining Software Repositories}
Mining Software Repositories (\textbf{MSR}) is a process of collecting and analysing big data from
repositories, which includes version control repositories, mailing list repositories, and bug
tracking repositories. MSR is applicable to a wide range of fields such as business, research, and
security \cite{poncin_2011}. The purpose of mining software repositories is to extract practical
information from rich metadata and discover hidden trends about a specific evolutionary
characteristic \cite{kagdi_2007}. The information collected could be used in various development
process. For example, some developers could gain insight by mining repositories, which may help them
to enhance their software quality based on previous implementation evidence of other developers
\cite{hassan_2008}. While MSR have various usages in different areas, the primary objective of this
project will be focusing on finding the security issues in open source software repositories through
MSR.

In order to identify both hidden and publicly disclosed patches, it is required to make effective
use of MSR technique. A MSR process is normally carried out using tools or scripts made by the
researchers themselves. Although there are many researches in the MSR field in the recent years, the
majority of the tools or scrips used are not published publicly \cite{robles_2010}. As a result, it
is not possible to fully replicate the previous research methods and make improvements based on
that. Despite the undisclosed information of research methods in many papers, Shang
\cite{shang_2009} suggested that the MSR process should be splitted into several stages, with each
stage focusing on specific topic of the problem to achieve the optimal efficiency.

\subsection{Keywords Search}
Keywords search is the core procedure of retrieving information from a repository. The prerequisite
is that the repository must have sufficient amount of valuable information, which can be estimated
by judging the history of the repository. To correctly and precisely retrieve the information from a
query, it is required to integrate some algorithms and modules into the search function.
% bookmark
Matsushita, Sasaki, and Inoue \cite{matsushita_2005} developed a repository search system
that make use of two functions: lexical analysis function and token comparing function. The system
produced very detailed results by deploying recursive search strategy into every commit. On the
contrary, Mockus and Votta \cite{mockus_2000} designed an automated program that

\subsection{Vulnerability Patch Prediction}
Williams and Hollingsworth \cite{williams_2005} developed a source code analysis tool that searches
for bug fixes and combines with information mined from repositories to improve the results. It is
stated that the most efficient way to utilise the historical information is to ignore the commit
messages and focus on mining the code changes. In order to locate the actual code changes for the
bug fix, a function return value checker was implemented to compares the number of warnings produced
by a same function across different versions. Williams and Hollingsworth assumed that a bug is fixed
if the warnings produced by the same function have decreased between two versions, and the final
result produced is a list of functions that are related to a potential bug fix in the commit
history.

\subsection{Vulnerabilities Classification Techniques}
Many classification techniques of computer attacks have been introduced before open source became
popular, and these early works have significant influence on the later work of finding and
classifying security issues in open source software repositories. An early work of Lindqvist and
Jonsson \cite{lindqvist_1999} implied that the \textit{location} of a computer flaw is a determining
factor in the classification process, in which their results are obtained through real penetration
attacks. Hansman and Hunt \cite{hansman_2005} adapted the ideas of Lindqvist and Jonsson and
suggested the concept of dimensions, where they categorised the attacks into several layers.

\subsection{Mining Vulnerabilities}
Cowan \cite{cowan_2003} suggested a list of software auditing and vulnerability mitigation tools
that perform static code and runtime analysis respectively to find bugs and vulnerabilities. This
approach is popular because it is able to find simple bugs in a short amount of time. However,
Bessey et al. \cite{bessey_2010} claimed that static tools have negative effect on technical
develoepment due to its high false positives nature. While this statement might be true, it does not
imply that all static tools are not effective as they differ in the techniques used in finding
vulnerabilities \cite{moser_2008}. It might require several experiments of different configurations
to obtain the best result, and the result may vary across different data sets.

This project extends prior work on Reis and Abreu's \cite{secbench} Secbench Mining Tool. The tool
aims to find vulnerabilities patch in GitHub repositories by using specific regular expressions for
each vulnerability pattern. Then it creates a test case for every vulnerability found and these test
cases are evaluated manually. Reis and Abreu \cite{reis_2017} discussed the procedure of the
evaluation and explained that human errors could occur due to source code complexity and similarity
of vulnerability pattern. The approach of Secbench Mining Tool is similar to the concept of this
project. However, it is not practical to perform manual evaluation on every result. In this project,
the tool developed should be able to automate the evaulation process to some extent, while
preserving the accuracy of the results.

\chapter{Requirements and Analysis}
The purpose of this chapter is to express the aims in more details, and discuss the problems to be
solved. This chapter will outline the requirements of the project and list the criteria to be met.
The analysis part will cover every aspect of the design, implementation and testing stage to ensure
that the project is feasible.

\section{Project Objectives}
Initially, the objectives set in \hyperref[sec:objectives]{\textbf{Section 1.2}} are an ideal
concept of this project. Having completed the background research and literature review, it is now
possible to provide a detailed description and more clearly defined objectives that improve the
feasibility of this project.

\begin{enumerate}
	\item \textbf{Vulnerability patterns}: The term `vulnerability pattern' is used to represent the
  commit message pattern of different vulnerabilities. Correctly identifying the regular
  expression of each vulnerability pattern is a time-comsuming process. It would need considerable
  refinement throughout the whole project. Hence, it could be more appropriate to reuse and
  improve the patterns provided in previous related work.
	\item \textbf{Mining the commits}: This task involves creating a repository mining tool that makes
	extensive use of the pre-defined regular expressions to search for the relative commits. It will
	be necessary to consider how closely a commit needs to match with the patterns for it to be
	included in the result. The file format of the result will be discussed in the later section.
	\item \textbf{Evaluating the mined commits}: The mining tool can be extended to include a separate
	function that evaluates the commits mined to find the actual code commit addressing the security
	vulnerabilities. However, all previous related work performed did not use automated techniques for
	the evaluation. This project will consider to automate the evaluation process to some extent while
	maintaining the accuracy of the results at the standard level.
\end{enumerate}

\section{Software Specification}
\begin{table}[H]
	\begin{center}
		\begin{tabular}{|p{10.3cm}|c|}
			\hline
			\rowcolor[HTML]{D8D8D8}
			\multicolumn{1}{|c|}{Criteria} & Importance \\ \hline
			\textbf{Compatibility}: The mining tool should be able to run on all machines that meet the
			system requirements. & Essential \\ \hline
			\textbf{Completeness}: The mining tool should be able to find all relevant commits of security
			vulnerabilities based on the regular expressions. & Essential \\ \hline
			\textbf{Repeatable}: The results should be repeatable and reproducible. & Essential \\ \hline
			\textbf{Robustness}: The mining tool should be able to handle all possible errors without
			terminating the mining process. & Essential \\ \hline
			\textbf{Scalability}: The mining tool should be able to work on different project sizes,
			provided that the repository contains a certain amount of information. & Essential \\ \hline
			\textbf{Automated Evaluation}: The process of classifying and evaluating the commits into
			different vulnerabilities patch should be automated to a certain extent. & Desirable \\ \hline
			\end{tabular}
		\captionof{table}{Specification of the mining tool} \label{table:criteria}
	\end{center}
\end{table}

\section{Considerations}
The aim of this section is to contemplate the options available for this project and review some of
the fundamental decisions to be made before the implementation.

\subsection{Programming Language}
Python 3 \cite{python} is chosen to be the main programming language for the repository mining tool.
While other programming languages may be more suitable for tackling specific problems of this
project, Python 3 provides sufficient coverage over every aspect with its comprehensive
functionality. The greatest advantage of Python 3 is that it has a wide range of libraries that
facilitate the development environment, which fully justified that a complete working solution can
be produced using Python 3.

\subsection{Libraries and Tools}
Since the mining tool is decided to be programmed in Python 3, a wide range of libraries could be
integrated to enhance its functionality.
\begin{itemize}
	\item PyGithub is a Python library build to access the GitHub API \cite{pygithub}.
	\item GitPython is a Python library build to interact with Git repositories using a combination of
	python and git command implementation \cite{gitpython}.
\end{itemize}

\subsection{File Format of Result}
The JavaScript Object Notation (\textbf{JSON}) \cite{json} has been chosen as the file format for
storing the results in this project. This is because JSON is supported in Python and it does not
require complicated operations in Python to access the data. While various alternative data
interchange formats such as the Extensible Markup Language (\textbf{XML}) \cite{xml} has its unique
advantages, it is important to choose a data interchange format that consumes less resource and have
lower processing time for large amount of data. Since it has been proved that JSON has better
performance than XML in terms of processing time and resource utilisation \cite{nurseitov_2009}, it
is considered that JSON would be the best option for this project.

\section{Proposed Method}
This project strongly emphasises the need for finding security issues in open source repositories
through MSR. While it might be impossible to discover the security patches in a repository through a
single search, the problem could be solved using divide and conquer. The ideal concept of this
project is to build a command-line interface program that is able to run two separate process: the
\textbf{mining} process and the \textbf{evaluation} process. The \textbf{mining} process takes a Git
repository as input, searches through the commit log, and stores the list of commits that might
potentially contain a patch in a JSON file. The \textbf{evaluation} process takes a JSON file as
input, and check the code difference of every commit in the log file to identify the real patches.

\section{Problems and Constraints}
As mentioned in \hyperref[sec:challenges]{\textbf{Section 1.3}}, the main challenges of this project
are \textbf{data}, \textbf{misclassification}, \textbf{evaluation} and \textbf{time}. This section
will discuss the problems in detail and review several ways of mitigating them, as well as analysing
the possible constraints that might affect the progress of the project.

It is expected that some problems might be solved and new problems could emerge in the course of the
project. There are several ways of mitigating the problems to reduce the risk, provided that the
problems are clearly identified and they are under the project scope. It is estimated that the
\textbf{evaluation} process would be the biggest challenge of this project since it was regarded as
a complicated and difficult area in previous researches. Moreover, this project plans to implement
an automated version of the evaluation process, which will further increases the difficulty level.

Automated evaulation is hard and does not guarantee to provide a good result. It is also extremely
challenging for the mining tool to work across repositories programmed in different programming
languages. The constraint is that the tool has to be exhaustively tested to find the optimal
threshold value and for it to be automated and produce good results. Although the tool might produce
good results on some repositories, it does not indicate that the tool will produce consistent
results on all repositories. To ensure the minimum quality of the results, one of the solutions
might be using both automated method for basic filtering and manual method for advance refinement.

\section{Testing}
This section covers a brief overview of the testing stage. It will be necessary to consider some of
the self-created test cases and scenarios in advance to find all possible bugs and flaws.

\subsection{Unit Testing}
Python provides an unit testing framework as part of its standard library, known as unittest
\cite{unittest}, which offers a complete set of functions suffice to cover the testing stage of this
project.

\subsection{System Testing}
After completing the unit testing, the mining tool has to be tested for its completeness and
robustness, as mentioned in \hyperref[table:criteria]{\textbf{Table 3.1}}. It is expected that the
program would not be able to handle complicated errors during the early implementation, and the
project schedule would become a iterative process between implementation and testing. It is assumed
that the testing stage would be the most time-consuming process in the whole project, thus it might
be required to allocate more time and effort into this stage.

\section{Evaluation}
This section briefly discusses the approach to evaluate the mining tool on the real world projects
to ensure that the requirements and criteria listed are practical and feasible.

\subsection{Real World Projects Evaluation}
The mining tool will be evaluated by testing it on real world open source projects. For the mining
tool to be beneficial to the public, it must be able to produce results with certain standard. This
could be validated by verifying the accuracy and relevance of the results. It is presumed that the
mining tool would only be suitable for a small set of repositories, and it might require to
experiment with different configurations to achieve the best result.

\section{Ethical Issues}
In this project, it has to be clearly declared that any known or unknown vulnerabilities found by
the mining tool in any repositories will not be publicly disclosed without the permission of the
original authors. The reason is that publishing the vulnerabilities publicly would make the
softwares highly vulnerable to attackers \cite{arora_2010}, and it is recommended to wait for the
official announcement from the software vendors.

\chapter*{Conclusions and Project Plan}
\section*{Plan of Action}
\subsection*{Semester 1 and Christmas Vacation}
\begin{ganttchart}[
		hgrid=true,
		vgrid={draw=none, dotted},
		expand chart=\textwidth
	]{1}{12}
	\gantttitle{Semester 1}{12} \\
	\gantttitlelist{1,...,12}{1} \\
	\ganttbar{Background Reading}{1}{10} \\
	\ganttbar{Introduction}{1}{1} \\
	\ganttbar{Analysis}{1}{1} \\
	\ganttbar{Literature Review}{2}{4} \\
	\ganttbar{Requirements and Analysis}{5}{6} \\
	\ganttbar{Proof reading and amendment}{7}{9} \\
	\ganttbar{Design and Implement}{10}{12}
	\ganttlink{elem2}{elem3}
	\ganttlink{elem3}{elem4}
	\ganttlink{elem4}{elem5}
	\ganttlink{elem5}{elem6}
\end{ganttchart}

\begin{itemize}
	\item \textbf{Week 7}: Starting from this week, discuss with the supervisor weekly about the
	document, also it is best to start the design stage early, and show the prototype to the
	supervisor.
\end{itemize}

\subsection*{Semester 2}
\begin{ganttchart}[
	hgrid=true,
	vgrid={draw=none, dotted},
	expand chart=\textwidth
]{1}{12}
\gantttitle{Semester 2}{12} \\
\gantttitlelist{1,...,12}{1} \\
\ganttbar{Design}{1}{1} \\
\ganttbar{Implement prototype}{2}{4} \\
\ganttbar{Testing prototype}{3}{6} \\
\ganttbar{Write dissertation}{7}{9} \\
\ganttbar{Submit project}{10}{10} \\
\ganttbar{Poster session}{11}{11}
\ganttlink{elem0}{elem1}
\ganttlink{elem2}{elem3}
\ganttlink{elem3}{elem4}
\ganttlink{elem4}{elem5}
\end{ganttchart}

\begin{itemize}
	\item \textbf{Week 1}: Finished the design and started implementation during holiday.
	\item \textbf{Week 2}: The implementing and testing stage is a repetitive process. It is very
	likely that the program will run into errors in the testing and had to spend more time fixing it.
\end{itemize}

% \chapter{Design}

% \chapter{Implementation and Testing}

% \chapter{Results and Discussion}

% \chapter{Conclusion}

\printbibliography[heading=bibintoc]

\end{document}